% Rapport TP3 - Agentic RAG sur la jurisprudence
% Compilation : pdflatex rapport.tex (deux fois), ou xelatex, ou Overleaf (dossier docs/rapport complet).
\documentclass[11pt,a4paper]{article}

\newcommand{\auteur}{Anas El Alaoui} % <-- à vérifier
\newcommand{\cours}{IA générative --- Séance 3 : Agents}

\usepackage{iftex}
\ifPDFTeX
  \usepackage[T1]{fontenc}
  \usepackage[utf8]{inputenc}
  \usepackage{lmodern}
  \usepackage{newunicodechar}
  \newunicodechar{→}{\ensuremath{\rightarrow}}
  \newunicodechar{≈}{\ensuremath{\approx}}
  \newunicodechar{✓}{\checkmark}
\else
  \usepackage{fontspec}
\fi
\usepackage[french]{babel}
\usepackage[a4paper,margin=2.2cm,headheight=14pt]{geometry}
\usepackage{microtype}
\usepackage{amssymb}
\usepackage{graphicx}
\usepackage{booktabs,longtable,tabularx,array}
\usepackage[table,dvipsnames]{xcolor}
\usepackage{tikz}
\usetikzlibrary{arrows.meta,positioning,shapes.geometric,fit,backgrounds,calc}
\usepackage{enumitem}
\usepackage{caption}
\usepackage{fancyhdr}
\usepackage{titlesec}
\usepackage[most]{tcolorbox}
\usepackage{listings}
\usepackage[hidelinks]{hyperref}
\usepackage{url}

% ---------- couleurs et mise en forme
\definecolor{navy}{HTML}{1F3A5F}
\definecolor{navy2}{HTML}{2D5A8E}
\definecolor{gold}{HTML}{A7782F}
\definecolor{soft}{HTML}{E8EEF6}
\definecolor{cDecision}{HTML}{2B5EA7}
\definecolor{cTool}{HTML}{6440A8}
\definecolor{cObs}{HTML}{6B7385}
\definecolor{cGen}{HTML}{1F3A5F}
\definecolor{cCheck}{HTML}{8A5A00}
\definecolor{cFinal}{HTML}{17734A}
\definecolor{cError}{HTML}{B3261E}
\newcommand{\tracetag}[2]{\textcolor{#1}{\footnotesize\bfseries\sffamily #2}}

\titleformat{\section}{\Large\bfseries\color{navy}}{\thesection}{0.8em}{}
\titleformat{\subsection}{\large\bfseries\color{navy2}}{\thesubsection}{0.7em}{}
\titleformat{\paragraph}[runin]{\normalsize\bfseries\color{navy}}{}{}{}[]
\setlength{\parskip}{0.45em}
\setlength{\parindent}{0pt}
\setlist{itemsep=0.15em,topsep=0.3em}
\captionsetup{font=small,labelfont={bf,color=navy}}

\pagestyle{fancy}
\fancyhf{}
\fancyhead[L]{\small\color{cObs}TP3 --- Agentic RAG sur la jurisprudence}
\fancyhead[R]{\small\color{cObs}\auteur}
\fancyfoot[C]{\small\thepage}
\renewcommand{\headrulewidth}{0.4pt}

\newtcolorbox{answerbox}[1]{enhanced,colback=soft,colframe=navy2,boxrule=0.5pt,arc=2mm,
  fonttitle=\bfseries\small,title={#1},left=2mm,right=2mm,top=1mm,bottom=1mm,breakable}
\newtcolorbox{keybox}{enhanced,colback=white,colframe=gold,boxrule=0.6pt,arc=2mm,left=3mm,right=3mm,
  borderline west={2.5pt}{0pt}{gold}}

\lstset{basicstyle=\ttfamily\small,columns=fullflexible,keepspaces=true,frame=single,
  rulecolor=\color{cObs!40},backgroundcolor=\color{soft!50},breaklines=true}

\newcommand{\shot}[3]{% \shot{fichier}{largeur}{légende}
  \IfFileExists{img/#1}{\begin{figure}[htbp]\centering
  \includegraphics[width=#2\textwidth]{img/#1}\caption{#3}\end{figure}}{}}
\newcommand{\gen}[1]{\IfFileExists{generated/#1}{\input{generated/#1}}{\emph{(fichier \texttt{\detokenize{generated/#1}} à générer avec \texttt{python -m scripts.make\_report\_tables})}}}

\begin{document}

% ======================================================================= titre
\begin{titlepage}
\begin{tikzpicture}[remember picture,overlay]
  \fill[navy] (current page.north west) rectangle ([yshift=-7.2cm]current page.north east);
  \fill[gold] ([yshift=-7.2cm]current page.north west) rectangle ([yshift=-7.45cm]current page.north east);
\end{tikzpicture}
\vspace*{0.4cm}
{\color{white}\sffamily
{\large \cours}\\[1.6cm]
{\fontsize{30}{36}\selectfont\bfseries Agentic RAG}\\[0.25cm]
{\Large sur la jurisprudence de la Cour de cassation}\par}
\vspace{3.2cm}
{\large\bfseries\color{navy} Travaux pratiques n°3}\\[0.2cm]
{\color{cObs} Du RAG statique (TP2) à un assistant juridique qui décide, appelle des outils,
vérifie ses preuves et trace son exécution, avec un LLM local.}

\vspace{1.4cm}
\begin{keybox}
\begin{tabularx}{\linewidth}{@{}l X@{}}
\textbf{Corpus} & 1\,973 décisions de la Cour de cassation (2015--2025, publiées au Bulletin), 26\,568 passages \\
\textbf{Index} & hybride BM25 + embeddings \texttt{multilingual-e5-small}, fusion RRF \\
\textbf{LLM} & \texttt{llama3.2:3b} servi par Ollama, en local sur CPU \\
\textbf{Agent} & LangGraph : agent $\rightarrow$ outils $\rightarrow$ réponse $\rightarrow$ vérificateur, 5 étapes max \\
\textbf{Outils} & \texttt{search\_cases}, \texttt{read\_case}, \texttt{external\_search} (Wikipédia FR) \\
\textbf{Interface} & application web locale, trace en direct (SSE), sources corpus / externes séparées \\
\end{tabularx}
\end{keybox}

\vfill
{\large \auteur}\hfill{\large \today}
\end{titlepage}

\tableofcontents
\newpage

% ======================================================================= 1
\section{Objectif et démarche}

Le TP2 a produit un RAG \emph{statique} : pour chaque question, la même chaîne s'exécute
(une recherche, des passages, une génération). Le TP3 transforme ce système en \textbf{agent} :
le modèle choisit l'action suivante (chercher, lire une affaire, consulter une source externe,
répondre), le code exécute l'outil, l'observation revient dans l'état, et un
\textbf{vérificateur} décide si les preuves suffisent avant de laisser passer la réponse.

Trois principes ont guidé la conception :
\begin{enumerate}
  \item \textbf{Réutiliser le TP2 tel quel comme outil.} Le corpus, l'index et la recherche ne
        changent pas : l'agent les appelle à travers \texttt{search\_cases} et \texttt{read\_case}.
  \item \textbf{Contrôler un petit modèle par le code.} Un modèle de 3\,B paramètres sur CPU fait des
        choix maladroits ; des \emph{garde-fous} explicites, visibles dans la trace, corrigent les
        décisions invalides au lieu de les cacher.
  \item \textbf{Tout rendre observable.} Chaque décision, appel d'outil, observation et vérification
        est publiée en direct dans l'interface et sauvegardée (JSON + Markdown).
\end{enumerate}

% ======================================================================= 2
\section{Point de départ : le RAG du TP2}

\paragraph{Données.} Le fichier \texttt{cour\_de\_cassation.jsonl.gz} du dataset
\texttt{antoinejeannot/jurisprudence} (environ 1\,Go, 553\,000 décisions) est lu en flux, sans être
stocké. On garde les décisions rendues depuis 2015, \textbf{publiées au Bulletin} (celles qui font
jurisprudence) et disposant d'un découpage en zones, puis on tire un échantillon équilibré par chambre
(tirage par réservoir) : 1\,973 décisions distinctes, dont 94\,\% ont un \emph{sommaire} officiel.

\paragraph{Découpage.} Le dataset indique pour chaque décision la position des zones
\emph{introduction, exposé du litige, moyens, motivations, dispositif}. Le découpage suit ces zones
(fenêtres de 170 mots, recouvrement de 35 mots) : un passage garde ainsi son \emph{rôle juridique},
ce qui sera essentiel pour l'agent. Le sommaire et les thèmes forment un passage à part, souvent le
meilleur point d'entrée pour une recherche par notion.

\paragraph{Index hybride.} BM25 (tokenisation française, mots vides, racinisation Snowball) capte les
termes exacts du droit et les numéros d'articles ; les embeddings \texttt{multilingual-e5-small}
capturent les reformulations en langage naturel. Les deux classements sont fusionnés par
\emph{Reciprocal Rank Fusion} ($k=60$), insensible aux échelles de score. Un numéro de pourvoi
(\texttt{23-20.161}) donne un accès direct à la décision.

\paragraph{Résumé et questions-réponses.} Le résumé suit la grille du juriste (faits, procédure,
question de droit, décision, portée) ; la question-réponse travaille sur les passages de l'affaire
sélectionnée. Dans les deux cas le modèle doit citer ses sources \texttt{[P$n$]}, en JSON, et le code
vérifie que chaque citation renvoie à un passage réellement fourni.

\shot{tp2_recherche.png}{0.98}{Interface du TP2 : recherche hybride, résultats avec le passage le plus pertinent, fiche de l'affaire et résumé structuré sourcé.}

% ======================================================================= 3
\section{Architecture agentique}

\begin{figure}[htbp]
\centering
\begin{tikzpicture}[
  box/.style={draw=navy2,rounded corners=3pt,fill=white,align=center,font=\small,minimum height=1.3cm,text width=2.7cm,line width=0.7pt},
  dec/.style={diamond,aspect=1.8,draw=gold,fill=gold!8,align=center,font=\small,inner sep=0pt,line width=0.7pt},
  term/.style={rounded corners=8pt,fill=navy,text=white,align=center,font=\small\bfseries,minimum height=0.9cm,text width=2.4cm},
  store/.style={cylinder,shape border rotate=90,aspect=0.2,draw=cObs,fill=soft,align=center,font=\footnotesize,minimum height=1.0cm,text width=2.2cm},
  arr/.style={-{Stealth[length=2.6mm]},line width=0.7pt,color=navy2},
  lab/.style={font=\scriptsize,color=cObs,fill=white,inner sep=1.5pt}]
  \node[term]  (q)      at (2.6, 2.0)  {Question};
  \node[box]   (agent)  at (2.6, 0)    {\textbf{agent}\\[1pt]\footnotesize LLM : action en JSON\\ + garde-fous en code};
  \node[box]   (tools)  at (8.4, 0)    {\textbf{tools}\\[1pt]\footnotesize exécute l'outil,\\ ajoute les preuves};
  \node[store] (idx)    at (12.6, 1.0) {Index TP2\\ BM25 + e5};
  \node[store] (ext)    at (12.6, -0.9) {API externe\\ Wikipédia FR};
  \node[box]   (answer) at (2.6, -3.0) {\textbf{answer}\\[1pt]\footnotesize rédige à partir\\ des preuves [C\#/E\#]};
  \node[dec]   (verify) at (8.4, -3.0) {\textbf{verify}\\[-1pt]\footnotesize preuves\\[-2pt]\footnotesize suffisantes ?};
  \node[term]  (final)  at (12.6, -3.0) {Final\\[-1pt]\footnotesize\mdseries réponse + sources};
  \draw[arr] (q) -- (agent);
  \draw[arr] ([yshift=3mm]agent.east) -- node[lab,above=1pt]{appel d'outil} ([yshift=3mm]tools.west);
  \draw[arr] ([yshift=-3mm]tools.west) -- node[lab,below=1pt]{observation} ([yshift=-3mm]agent.east);
  \draw[arr] (agent) -- node[lab]{answer} (answer);
  \draw[arr] (answer) -- (verify);
  \draw[arr] (verify) -- node[lab,above=2pt]{oui / limite} (final);
  \draw[arr,dashed,color=gold] (verify.north) -- ++(0,0.75) -| node[lab,pos=0.25,above=1pt]{non : motif de l'échec} ([xshift=8mm]agent.south);
  \draw[arr,dashed,color=cError] (agent.west) -- ++(-0.7,0) |- node[lab,pos=0.75,below=1pt]{limite atteinte sans preuve / erreur : arrêt contrôlé} ([yshift=-1.55cm]verify.south) -| (final.south);
  \draw[dotted,cObs,line width=0.9pt] (tools.east) -- (idx.west);
  \draw[dotted,cObs,line width=0.9pt] (tools.east) -- (ext.west);
\end{tikzpicture}
\caption{Graphe LangGraph de l'agent. Les flèches pleines sont le chemin nominal ; le vérificateur renvoie
vers l'agent avec le motif de l'échec tant que la limite d'étapes n'est pas atteinte.}
\label{fig:archi}
\end{figure}

\subsection{État partagé}
L'état (\texttt{AgentState}) ne contient que ce qui sert à décider :
\begin{center}\small
\begin{tabularx}{\textwidth}{@{}l X@{}}
\toprule
Champ & Contenu \\ \midrule
\texttt{question}, \texttt{memory} & question courante ; affaire et échange précédents (bonus mémoire) \\
\texttt{evidence} & preuves citables numérotées : \texttt{C1, C2\ldots} (corpus), \texttt{E1\ldots} (externe) \\
\texttt{candidates} & décisions trouvées, étiquetées \texttt{D1, D2\ldots} (l'agent manipule ces étiquettes courtes plutôt que des identifiants longs) \\
\texttt{searches}, \texttt{read\_cases} & historique des recherches et des affaires lues (détection des répétitions) \\
\texttt{iteration}, \texttt{max\_iterations} & compteur de décisions et limite (5 par défaut) \\
\texttt{feedback}, \texttt{verification} & motif du dernier échec de vérification, détail des contrôles \\
\texttt{draft}, \texttt{final}, \texttt{status} & brouillon, réponse finale avec sources, statut (vérifiée, non vérifiée, arrêt, erreur) \\
\bottomrule
\end{tabularx}
\end{center}
La trace, les compteurs (appels LLM, tokens, latence) et la validation humaine vivent dans un
\texttt{RunContext} transmis par la configuration du graphe : ils observent l'exécution sans polluer
l'état métier.

\subsection{Outils}
\begin{center}\small
\begin{tabularx}{\textwidth}{@{}l X@{}}
\toprule
Outil & Comportement \\ \midrule
\texttt{search\_cases(query, k)} & recherche hybride du TP2, regroupée par décision ; renvoie titre, solution et extrait. Les trois premiers résultats fournissent déjà une preuve citable (souvent le sommaire). \\
\texttt{read\_case(case\_id)} & lecture « juridique » : sommaire officiel + le passage le plus proche de la question dans chacune des zones \emph{exposé, moyens, motivations, dispositif}. La question est enrichie du vocabulaire des arrêts (« fait grief », « rejette »\ldots). \\
\texttt{external\_search(query)} & API MediaWiki de Wikipédia FR (recherche + résumé introductif). Les résultats sont marqués \textbf{externes} et ne peuvent jamais être présentés comme de la jurisprudence. \\
\bottomrule
\end{tabularx}
\end{center}
Une erreur d'outil (\texttt{ToolError} : décision introuvable, API indisponible, requête vide) est
affichée dans la trace en rouge, enregistrée dans l'état, et l'agent continue ; une erreur du LLM
provoque un arrêt contrôlé.

\subsection{Décision de l'agent et garde-fous}
À chaque tour, le LLM reçoit un état résumé (question, recherches faites, décisions trouvées avec leurs
étiquettes, affaires déjà lues, nombre de preuves, retour du vérificateur, étape $i$/5) et répond en JSON :
\texttt{\{"reason", "action", "args"\}}. La justification est une phrase courte : la trace est
\emph{opérationnelle}, pas un long raisonnement interne.

Le modèle 3\,B se trompe de façon prévisible ; le code applique alors des garde-fous, chacun affiché
dans la trace :
\begin{itemize}
  \item action inconnue $\rightarrow$ recherche (ou réponse si des preuves existent) ;
  \item réponse demandée sans aucune preuve $\rightarrow$ recherche ;
  \item recherche identique à une précédente $\rightarrow$ réponse si une affaire a été lue, sinon lecture du meilleur candidat, sinon source externe, sinon arrêt ;
  \item affaire déjà lue $\rightarrow$ réponse (ou nouvelle affaire si le vérificateur a demandé plus de preuves) ;
  \item numéro de pourvoi cité dans la question et décision trouvée mais non lue $\rightarrow$ lecture obligatoire ;
  \item nouvelle réponse après un refus du vérificateur sans nouvelle preuve $\rightarrow$ nouvelle action de recherche.
\end{itemize}

\subsection{Vérification avant la réponse finale}
Le nœud \texttt{verify} applique des contrôles déterministes puis un juge LLM :
\begin{center}\small
\begin{tabularx}{\textwidth}{@{}l X@{}}
\toprule
Contrôle & Règle \\ \midrule
Preuves présentes & au moins un passage du corpus dans l'état \\
Citations valides & la réponse cite au moins une preuve, et uniquement des preuves existantes \\
Décision du corpus citée & une réponse appuyée seulement sur Wikipédia est refusée \\
Affaire demandée citée & si la question vise un numéro de pourvoi, les passages cités viennent de cette affaire \\
Notions clés sourcées & un terme rare de la question (IDF $\geq 6$, ex. « imprévision ») repris dans la réponse doit figurer dans les preuves \\
Pas d'information non sourcée & chaque phrase partage au moins 40\,\% de ses mots (racinisés) avec les preuves fournies \\
Cohérence réponse / passages & juge LLM (« chaque affirmation est-elle justifiée par les passages ? ») ; une objection n'est retenue que si l'affirmation est réellement peu présente dans les preuves \\
\bottomrule
\end{tabularx}
\end{center}
En cas d'échec, le motif (« la réponse s'appuie sur une autre affaire », « phrase non soutenue »,
« manque : \ldots ») devient le \texttt{feedback} de l'agent, qui doit relancer une recherche,
reformuler ou lire une autre affaire. À la limite d'étapes, la réponse est rendue avec le statut
\emph{non vérifiée} et un avertissement visible.

\subsection{Interface et exécution}
Le run tourne dans un thread côté serveur (\texttt{graph.stream(..., stream\_mode="updates")}) ;
chaque nœud publie ses événements dans une file, diffusée au navigateur en \emph{Server-Sent Events}.
L'utilisateur voit en direct le nœud actif du graphe, l'outil appelé avec ses arguments, le résumé de
l'observation, le résultat de la vérification et la condition d'arrêt. La réponse finale sépare les
\textbf{sources du corpus} et les \textbf{sources externes} ; les citations sont cliquables.

\shot{tp3_agent.png}{0.98}{Onglet « Agent juridique » : options (source externe, validation humaine, juge LLM, mémoire, limite d'étapes), graphe avec le nœud actif, trace étape par étape et réponse avec ses sources et ses mesures.}

\subsection{Bonus réalisés}
\begin{itemize}
  \item \textbf{Mémoire} : l'affaire citée dans la réponse précédente est proposée à l'agent sous l'étiquette \texttt{D0} (question de suivi : « et quelle a été la solution ? »).
  \item \textbf{Human-in-the-loop} : option de validation avant tout appel externe ; un refus est enregistré et l'agent continue sans la source externe.
  \item \textbf{Observabilité} : appels LLM, tokens, temps passé dans le modèle et durée totale, par run et par appel.
  \item \textbf{Rôles séparés} : le rédacteur et le vérificateur sont deux nœuds distincts avec des consignes différentes ; le superviseur (le graphe) décide quand appeler l'un ou l'autre.
\end{itemize}

% ======================================================================= 4
\section{Tests des trois scénarios}

Les traces ci-dessous sont générées automatiquement à partir des exécutions réelles
(\texttt{python -m scripts.run\_scenarios}, fichiers dans \texttt{docs/traces/}). Les temps sont mesurés
sur CPU, sans GPU.

\subsection{Scénario 1 --- RAG simple}
\emph{Question : Le juge judiciaire peut-il apprécier le caractère réel et sérieux du licenciement d'un
salarié protégé autorisé par l'inspecteur du travail ?} Une seule recherche doit suffire : le sommaire
officiel d'un arrêt de la chambre sociale énonce la règle. On vérifie que l'agent ne multiplie pas les
étapes.

{\small\gen{trace_s1.tex}}

\subsection{Scénario 2 --- Recherche multi-étapes}
\emph{Question : Dans l'affaire n° 23-20.161 sur un bail rural, que reprochait le bailleur à l'arrêt d'appel
et qu'a décidé la Cour de cassation ?} Il faut trouver la décision puis lire précisément le moyen
(« fait grief à l'arrêt\ldots ») et la réponse de la Cour.

{\small\gen{trace_s2.tex}}

\subsection{Scénario 3 --- Information insuffisante}
\emph{Question : Qu'est-ce que la théorie de l'imprévision et la Cour de cassation l'a-t-elle appliquée à un
contrat ?} La notion est générale et le corpus ne contient pas forcément d'arrêt qui la nomme : le premier
retrieval est insuffisant, l'agent doit reformuler ou consulter la source externe, sans présenter
Wikipédia comme de la jurisprudence.

{\small\gen{trace_s3.tex}}

% ======================================================================= 5
\section{Comparaison avec le RAG statique du TP2}

Le RAG statique exécute, pour la même question, une recherche hybride, prend 5 passages (2 au plus par
décision) et génère une réponse citée en un seul appel LLM.

\begin{table}[htbp]
\centering\small
\gen{comparaison.tex}
\caption{Coût et résultat de l'agent et du RAG statique sur les trois scénarios (même modèle, même index).}
\end{table}

{\small\gen{comparaison_reponses.tex}}

\paragraph{Où l'agent apporte une vraie valeur.} Sur le scénario 2, la question vise une affaire
précise et deux éléments de la décision (le grief du pourvoi et la solution). Le RAG statique ne répond
qu'à moitié : il recopie la règle posée par la Cour, sans le grief ni la solution, et sans citer de passage.
L'agent trouve la décision par son numéro, la \emph{lit} zone par zone (moyen, motivations, dispositif) et
produit une réponse complète, citée et vérifiée, en 64\,s. Sur le scénario 3, les deux systèmes
reprennent la même affirmation fausse (« la Cour a appliqué l'imprévision au contrat de travail ») ; mais
le RAG statique la présente comme un fait, alors que l'agent relance la recherche, consulte la source
externe puis rend sa réponse avec le statut \emph{non vérifiée} et le motif. Sur une question juridique,
savoir signaler qu'on ne sait pas vaut presque autant qu'une bonne réponse.

\paragraph{Où le RAG simple est préférable.} Pour une question de principe dont la réponse tient dans un
sommaire (scénario 1), le RAG statique donne la bonne conclusion en un seul appel LLM, en 15\,s contre
85\,s pour l'agent (4 appels LLM, environ trois fois plus de tokens). Sur CPU, chaque appel au modèle
3\,B coûte 10 à 40 secondes : la latence est le premier prix de l'agentivité. Le RAG statique est aussi
plus prévisible (toujours le même chemin) et plus simple à tester. Sa faiblesse est ailleurs : dans nos
trois essais, il n'a cité \emph{aucun} passage et a ajouté une justification inventée au scénario 1
(« il ne s'agit pas d'un licenciement économique »), ce qu'aucun contrôle ne détecte dans un pipeline fixe.

% ======================================================================= 6
\section{Analyse}

\subsection{Ce qui fonctionne}
\begin{itemize}
  \item Le chemin nominal \emph{recherche $\rightarrow$ lecture $\rightarrow$ réponse $\rightarrow$ vérification} est suivi sans étape inutile quand la question est claire.
  \item Les garde-fous rendent le système robuste aux erreurs d'un petit modèle, et ces corrections sont visibles dans la trace au lieu d'être silencieuses.
  \item La lecture « par zone » (exposé, moyens, motivations, dispositif) donne au rédacteur exactement les passages qu'un juriste lirait.
  \item Le vérificateur détecte réellement des erreurs : réponses fondées sur la mauvaise affaire, affirmations absentes des passages, réponses sans citation.
\end{itemize}

\subsection{Échecs observés}
\begin{itemize}
  \item \textbf{Boucle de recherche.} Dans les premières versions, le modèle 3\,B relançait la même recherche ou relisait l'affaire déjà lue ; un premier garde-fou le redirigeait vers des affaires sans rapport, et la réponse finale portait sur une autre décision. Correction : une affaire déjà lue mène à la réponse, et la sélection des preuves privilégie l'affaire citée par son numéro.
  \item \textbf{Vérification trompée par des preuves hors sujet.} Une réponse sur les « fumures et arrière-fumures » (affaire 23-19.267) a été validée pour une question sur l'affaire 23-20.161 : elle était bien soutenue par les passages cités, mais pas par la bonne affaire. D'où le contrôle « affaire demandée citée ».
  \item \textbf{Faux négatifs du juge LLM.} Le juge 3\,B a déclaré « non justifiée » une phrase recopiant presque mot pour mot le moyen du pourvoi, ce qui a déclenché des recherches inutiles jusqu'à la limite d'étapes. Ses objections sont désormais confrontées au contrôle lexical.
  \item \textbf{Vocabulaire.} « Que reprochait le bailleur ? » ne partage aucun mot avec « M. [T] fait grief à l'arrêt » : sans l'enrichissement de la requête par le vocabulaire des arrêts, le passage décisif n'était pas retrouvé.
  \item \textbf{Hallucination validée, puis corrigée.} Dans une version intermédiaire, la réponse du scénario 3 (« la Cour a appliqué l'imprévision au contrat de travail ») a été \emph{vérifiée} : le juge LLM l'avait bien rejetée, mais son objection avait été écartée parce que la plupart des mots de la phrase figuraient dans les preuves. Le mot décisif, « imprévision », n'y figurait pas. D'où le contrôle « notions clés sourcées » : un terme rare de la question (IDF $\geq 6$) repris dans la réponse doit apparaître dans les preuves.
  \item \textbf{Le rédacteur ignore la source externe.} Dans la trace finale du scénario 3, l'agent a bien récupéré la définition de Wikipédia, mais le modèle 3\,B rédige encore la même affirmation fausse au lieu de s'appuyer sur [E1]. Le vérificateur l'arrête : la réponse sort avec le statut \emph{non vérifiée}. Le système échoue proprement, mais il échoue.
\end{itemize}

\subsection{Améliorations possibles}
\begin{itemize}
  \item Un modèle plus grand ou spécialisé (7\,B quantifié, ou GPU) pour le juge et l'agent : la plupart des garde-fous compensent des erreurs de raisonnement du 3\,B.
  \item Un \emph{reranker} (cross-encoder multilingue) après l'index hybride, pour mieux ordonner les passages avant la lecture.
  \item Une vraie API juridique (Judilibre, Légifrance via PISTE) à la place de Wikipédia, pour citer les textes de loi visés.
  \item Une décision « règle d'abord, LLM ensuite » : router sans LLM les questions qui contiennent un numéro de pourvoi, et réserver le LLM aux cas ambigus.
  \item Autoriser une réponse \emph{négative} vérifiable (« aucune décision du corpus consulté n'applique cette notion », appuyée sur la liste des recherches faites), plutôt que d'exiger une décision du corpus citée : c'est la bonne réponse au scénario 3.
\end{itemize}

\begin{keybox}
\textbf{Question directrice : de meilleures décisions, ou seulement plus d'appels LLM ?}
Les deux, selon la question. Sur une question de principe, l'agent ajoute surtout des appels autour du
même RAG. Sur une question qui exige de choisir une source (une affaire précise, une notion absente du
corpus), il prend de meilleures décisions que le pipeline fixe, \emph{à condition} d'être encadré : avec
un modèle de 3\,B, la qualité vient autant des garde-fous et du vérificateur en code que du LLM lui-même.
\end{keybox}

% ======================================================================= annexe
\appendix
\section{Reproduire les résultats}
\begin{lstlisting}
pip install -r requirements.txt
ollama pull llama3.2:3b
python -m scripts.prepare_juris --n 2000      # corpus
python -m scripts.build_index                 # index BM25 + embeddings
uvicorn backend.main:app --port 8000          # interface : http://localhost:8000
python -m scripts.run_scenarios               # 3 scénarios + RAG statique -> docs/traces/
python -m scripts.make_report_tables          # tableaux LaTeX de ce rapport
pytest -q                                     # tests (faux LLM)
\end{lstlisting}

Organisation du code : \texttt{backend/juris/} (corpus, index, RAG du TP2), \texttt{backend/agent/}
(outils, graphe LangGraph, exécution et traces), \texttt{backend/juris\_api.py} (API),
\texttt{frontend/app/} (interface), \texttt{scripts/} (préparation, index, scénarios, rapport),
\texttt{tests/} (TP1 et agent, avec un faux LLM).

\end{document}
